\documentclass[conference]{IEEEtran}

% =========================================================
% PACKAGES
% =========================================================
\usepackage[utf8]{inputenc}
\usepackage[T1]{fontenc}
\usepackage{microtype}                % better justification and spacing
\usepackage{cite}
\usepackage{amsmath,amssymb}
\usepackage{graphicx}
\usepackage{tikz}
\usetikzlibrary{positioning,fit,backgrounds,arrows.meta}
\usepackage[table]{xcolor}
\usepackage{booktabs}
\usepackage{tabularx}
\usepackage{array}
\usepackage{enumitem}
\usepackage{listings}
\usepackage{balance}                  % balances the last page's columns
\usepackage{url}
\usepackage[colorlinks=true,linkcolor=navy,citecolor=navy,urlcolor=navy]{hyperref}

% -------------------- COLORS --------------------
\definecolor{navy}{RGB}{25,55,95}
\definecolor{rowgray}{gray}{0.94}
\definecolor{headgray}{gray}{0.86}

% -------------------- LISTS --------------------
\setlist{leftmargin=*, topsep=3pt, partopsep=0pt, itemsep=1.5pt, parsep=0pt}
\setlist[enumerate,1]{label=\arabic*.}

% -------------------- URL BREAKING --------------------
\makeatletter
\g@addto@macro\UrlBreaks{\do\/\do\-\do\_\do\.}
\makeatother
\Urlmuskip=0mu plus 1mu

% -------------------- TABLES --------------------
\newcolumntype{Y}{>{\raggedright\arraybackslash}X}
\newcolumntype{L}[1]{>{\raggedright\arraybackslash}p{#1}}
\newcolumntype{C}[1]{>{\centering\arraybackslash}p{#1}}
\renewcommand{\arraystretch}{1.18}
\setlength{\tabcolsep}{4pt}

% -------------------- CODE BLOCKS --------------------
\lstset{
  basicstyle=\ttfamily\footnotesize,
  backgroundcolor=\color{black!4},
  frame=single,
  rulecolor=\color{black!20},
  breaklines=true,
  columns=fullflexible,
  keepspaces=true,
  showstringspaces=false,
  xleftmargin=2pt,
  xrightmargin=2pt,
  aboveskip=5pt,
  belowskip=5pt
}

% -------------------- GENERAL SPACING --------------------
\setlength{\emergencystretch}{1.5em}
\setlength{\textfloatsep}{10pt plus 2pt minus 2pt}
\setlength{\floatsep}{8pt plus 2pt minus 2pt}
\setlength{\intextsep}{8pt plus 2pt minus 2pt}
\setlength{\abovecaptionskip}{4pt}
\setlength{\belowcaptionskip}{0pt}
\setlength{\abovedisplayskip}{4pt plus 1pt minus 1pt}
\setlength{\belowdisplayskip}{4pt plus 1pt minus 1pt}

\begin{document}

% =========================================================
% TITLE
% =========================================================
\title{Code Checking Assistant: An Automated Code Grading System}

\author{
\IEEEauthorblockN{Gaurav Chandel}
\IEEEauthorblockA{\textit{Roll No.\ 2611MC17}\\
\textit{Dept.\ of Computer Science \& Engineering}\\
\textit{Indian Institute of Technology Patna}}
\and
\IEEEauthorblockN{Prajodeep Patil}
\IEEEauthorblockA{\textit{Roll No.\ 2611MC11}\\
\textit{Dept.\ of Computer Science \& Engineering}\\
\textit{Indian Institute of Technology Patna}}
\and
\IEEEauthorblockN{Aditya Sharma}
\IEEEauthorblockA{\textit{Roll No.\ 2611CS11}\\
\textit{Dept.\ of Computer Science \& Engineering}\\
\textit{Indian Institute of Technology Patna}}
\and
\IEEEauthorblockN{Anshul Kaushal}
\IEEEauthorblockA{\textit{Roll No.\ 2611AI10}\\
\textit{Dept.\ of Computer Science \& Engineering}\\
\textit{Indian Institute of Technology Patna}}
}

\maketitle

% =========================================================
% ABSTRACT
% =========================================================
\begin{abstract}
The Code Checking Assistant is a Python system for automated assessment of C
programming submissions. An instructor supplies a \texttt{rules.json} file with
the assignment statement, a problem contract, and weighted metrics (nine by
default) worth 100 points in total. The submission is compiled and executed in a locked-down Docker
sandbox, where independent analyzers collect evidence: GCC diagnostics,
\texttt{cppcheck}, input/output and edge-case tests, Valgrind Memcheck,
AddressSanitizer/UndefinedBehaviorSanitizer, runtime profiles, syntax-tree facts
from \texttt{pycparser}, source metrics, and token-shingle similarity. When the
instructor's tests are not enough, the system generates more test inputs and
computes their expected outputs by running a reference program in the sandbox:
the instructor's solution when one is given, otherwise a program the model writes
from the assignment and that must reproduce the instructor's cases. A local Qwen3.5-2B model served by \texttt{llama.cpp} describes the
program, runs up to three recursive self-improvement (RSI) review passes that
must cite a deterministic evidence ledger, and rates (0--10) only the metrics
that no executed test or analyzer measured. Python converts the ratings to points
and computes the grade. Executed test pass rates and completed analyzers are
authoritative, and a metric without evidence is reported as unmeasured instead of
being given a neutral or zero score. Requests run at temperature 0 with a fixed
seed and no prompt cache, and generated test suites are frozen by fingerprint.
On a stratified 600-program FormAI-v2 sample, the static-tool detector reaches
63.3\% recall, 95.7\% specificity, and 93.6\% precision, and we report smaller
LLM-assisted runs with confidence intervals. The source code is available at
\url{https://github.com/Enthusiast101/Code-Grader}.
\end{abstract}

\begin{IEEEkeywords}
automated grading, C programming, Docker sandbox, Valgrind, sanitizers, local LLM,
test generation, FormAI dataset, vulnerability detection
\end{IEEEkeywords}

% =========================================================
% I. INTRODUCTION
% =========================================================
\section{Introduction}

\subsection{Background}
Programming assignments are commonly evaluated by compiling programs, executing
test cases, checking coding requirements, and assessing solution quality. Manual
evaluation takes considerable time, particularly for large numbers of
submissions. The Code Checking Assistant automates much of this process. It
accepts a student C program and instructor-defined grading rules, runs the
configured checks, calculates a score out of 100, and writes a structured JSON
report and a student-facing text report.

The approach combines conventional program analysis with a locally hosted
language model. Conventional tools compile, execute, and analyze the program, and
their measurements decide every metric they can measure. The language model
explains the evidence, finds defects that the tools do not report, and rates the
qualitative metrics (such as readability) that no tool measures directly.

\subsection{Problem Statement}
Manual checking of C assignments is time-consuming and can lead to inconsistent
feedback. A useful automated system should identify compilation errors, detect
execution and memory errors, evaluate the configured requirements, and give the
same score each time the same submission is graded. The main challenges are to
automate these tasks without letting student programs affect the host machine,
and to use a language model without letting unsupported model claims change the
grade.

\subsection{Project Objectives}
The project builds an automated C-code grader that:
\begin{enumerate}
    \item accepts a student C source file and an instructor \texttt{rules.json} with weighted metrics, limits, and test files;
    \item compiles and executes the submission inside a restricted Docker sandbox;
    \item checks correctness with instructor tests and with generated tests whose expected outputs come from executing a reference program;
    \item collects compiler, static-analysis, sanitizer, Valgrind, runtime, structural, and similarity evidence;
    \item records every fact in an evidence ledger with stable IDs, and keeps raw tool logs in the report;
    \item uses a local LLM for program understanding, a bounded RSI defect review, and rubric ratings of metrics that have no measured score;
    \item calculates the final grade in Python, with measured results taking precedence over model ratings;
    \item produces reproducible JSON and text reports; and
    \item measures vulnerability detection against the FormAI-v2 dataset.
\end{enumerate}

% =========================================================
% II. ARCHITECTURE
% =========================================================
\section{System Architecture}

\subsection{Overview}
The system is organized around a Python coordinator
(\path{grading_agent/main.py}) and independent analyzer modules registered in an
analyzer registry. A grading run proceeds as follows:
\begin{enumerate}
    \item \textbf{Rule loading:} \path{rule_loader.py} validates \texttt{rules.json}: assignment statement, problem contract, time and memory limits, metric names, analyzers, weights, point caps, and test-file paths.
    \item \textbf{Sandbox and preflight:} a Docker container is started, tool versions are recorded, and the submission is compiled once and smoke-run on one test. A compile failure or crash marks execution metrics as inconclusive.
    \item \textbf{Test suite:} instructor test files are loaded, and generated tests are added when enabled (Section~\ref{sec:testgen}).
    \item \textbf{Analyzers:} each enabled metric's analyzer runs independently. An analyzer failure becomes an \texttt{unavailable} record for that metric and does not stop the other analyzers.
    \item \textbf{Evidence ledger:} analyzer output, syntax-tree facts, and source lines become ledger entries with stable IDs.
    \item \textbf{LLM review:} program understanding, up to three RSI passes, deterministic prevalidation, model adjudication, and rubric ratings.
    \item \textbf{Scoring and reporting:} Python picks each metric's score by evidence priority, sums the points, and writes the JSON and text reports.
\end{enumerate}
A Gradio interface (\path{ui/app.py}) wraps this pipeline in two tabs:
\emph{Grade a C submission}, with separate \emph{Generate Test Cases} and
\emph{Grade Submission} actions, and \emph{FormAI benchmark}. A command-line
entry point (\path{grade.py}) runs the same pipeline.

\subsection{Technologies Used}
Table~\ref{tab:tech} lists the technology stack.

\begin{table}[!t]
\centering
\caption{Technology Stack}
\label{tab:tech}
\footnotesize
\rowcolors{2}{rowgray}{white}
\begin{tabularx}{\columnwidth}{@{}>{\bfseries\raggedright\arraybackslash}p{1.95cm} Y@{}}
\toprule
\rowcolor{headgray}
\textbf{Technology} & \textbf{Purpose} \\
\midrule
Python 3.12 & Coordinator, analyzers, scoring, reporting, benchmark; managed with \texttt{uv}. \\
Docker & Sandbox image \texttt{agent-sandbox:latest} (Debian slim with GCC, \texttt{cppcheck}, Valgrind), run as an unprivileged user. \\
GCC & C11 compilation with \texttt{-Wall -Wextra -Wpedantic}; sanitizer builds. \\
Valgrind Memcheck & Per-test leak and memory-error detection. \\
ASan/UBSan & Per-test memory and undefined-behaviour detection. \\
\texttt{cppcheck} & Static diagnostics for the structure metric and the benchmark. \\
\texttt{pycparser} & Syntax-tree facts and safety-property checks. \\
\texttt{lizard} & Cyclomatic complexity and function length (lexical fallback when unavailable). \\
\texttt{llama.cpp} & \texttt{llama-server} (build b12005 or later) with JSON-schema constrained decoding. \\
Qwen3.5-2B & Default local model, a 4-bit GGUF build. \\
Gradio & Web interface for grading and benchmarking. \\
SQLite & Per-program index of the FormAI datasets. \\
\bottomrule
\end{tabularx}
\end{table}

\subsection{External Reference Dataset: FormAI}
FormAI \cite{formai,formaiv2,formairepo} contains AI-generated C programs labelled
with the ESBMC bounded model checker \cite{esbmc}. The benchmark mainly uses
FormAI-v2 (programs written by nine LLMs, verified with ESBMC~7.5); FormAI-v1
(GPT-3.5 programs, ESBMC~7.3) is also supported. The published files contain one
record per verified property, so \path{formai_dataset.py} downloads the archives,
verifies their SHA-256 checksums, and builds a SQLite index with one row per
program. A program is \textsc{vulnerable} if any record has an ESBMC
counterexample, \textsc{safe} if every property was verified within the bound, and
\textsc{unknown} otherwise. Time-outs and parse errors are not evidence of safety,
so \textsc{unknown} programs are excluded. Of the 310,487 FormAI-v2 programs,
192,713 are vulnerable, 25,674 safe, and 92,100 unknown, which leaves 218,387
eligible programs. Each vulnerable line is mapped to one of nine families
(input buffer overflow, NULL dereference, out of bounds, invalid pointer, memory
leak, integer overflow, float overflow, division by zero, other).

The archives are stored under \path{resources/} and are not committed. FormAI is
used only for the security benchmark. It is not a grading-rule file and does not
replace assignment test cases. Building the index never compiles or runs a FormAI
program.

\subsection{System Architecture and Project Flowchart}
Fig.~\ref{fig:flow} shows the workflow. In grading mode, the instructor's rules,
the student program, and the test cases go through the sandbox, the analyzers, the
LLM review, and deterministic scoring. In benchmark mode, sampled FormAI sources go
through the source-level detectors, and the labels are read only after every
prediction has been written.

\begin{figure*}[!t]
\centering
\colorlet{detfill}{navy!8}
\colorlet{llmfill}{orange!18}
\colorlet{llmdraw}{orange!70!black}
\colorlet{datafill}{black!6}
\resizebox{\textwidth}{!}{%
\begin{tikzpicture}[
  font=\scriptsize,
  >={Stealth[length=4pt]},
  box/.style={draw=navy, rounded corners=2pt, align=center, text width=2.5cm,
              minimum height=1cm, inner sep=3pt, fill=detfill},
  llm/.style={box, fill=llmfill, draw=llmdraw},
  data/.style={box, fill=datafill, draw=black!50},
  flow/.style={->, semithick, draw=black!75},
  aux/.style={->, semithick, dashed, draw=black!55},
  lane/.style={draw=black!30, rounded corners=4pt, inner sep=6pt, fill=black!2},
  lanelabel/.style={font=\footnotesize\bfseries, text=navy, anchor=north west, inner sep=3pt},
]
% ---------------- Grading mode, row 1 (left to right) ----------------
\node[data] (in) at (0,0)
  {\textbf{Instructor inputs}\\ \texttt{rules.json}, \texttt{submission.c}, test files, optional reference solution};
\node[box] (rules) at (3.3,0)
  {\textbf{Rule loader}\\ validates contract, limits, analyzer parameters, 100-point caps};
\node[box] (sand) at (6.6,0)
  {\textbf{Docker sandbox}\\ no network, read-only root, CPU/memory/PID limits\\[1pt]
   \textbf{Preflight}: compile + smoke run};
\node[llm] (suite) at (9.9,0)
  {\textbf{Test suite}\\ instructor tests + generated inputs; outputs from a reference program; frozen by fingerprint};
\node[box, text width=3.4cm] (an) at (13.7,0)
  {\textbf{Analyzers (one per metric)}\\[1pt]
   \emph{Execution:} I/O tests, edge cases, Valgrind, ASan/UBSan, runtime\\[1pt]
   \emph{Source:} GCC, \texttt{cppcheck}, \texttt{-fanalyzer}, syntax tree, modularity, complexity, readability, similarity};
% ---------------- Grading mode, row 2 (right to left) ----------------
\node[box, text width=3.4cm] (led) at (13.7,-3.1)
  {\textbf{Evidence ledger}\\ stable IDs for test results, tool findings, syntax-tree facts, source lines};
\node[llm, text width=3.1cm, align=left] (rev) at (9.9,-3.1)
  {\textbf{LLM review}\\
   1.~program understanding\\
   2.~RSI, at most 3 passes\\
   3.~Python prevalidation\\
   4.~adjudication\\
   5.~0--10 levels for unmeasured metrics};
\node[box] (score) at (6.3,-3.1)
  {\textbf{Deterministic scoring}\\ tests and analyzers first, LLM levels for the rest; unmeasured $\Rightarrow$ no grade};
\node[data] (rep) at (2.9,-3.1)
  {\textbf{Reports}\\ JSON + text: points, provenance, findings, hashes, timings};
% ---------------- Shared local model ----------------
\node[llm, text width=3.1cm] (model) at (9.9,-5.6)
  {\textbf{Local LLM}\\ Qwen3.5-2B on \texttt{llama-server}; temperature 0, seed 42, JSON schema};
% ---------------- Benchmark mode ----------------
\node[data] (fa) at (0,-8.2)
  {\textbf{FormAI-v2}\\ archives, SHA-256 verified};
\node[box] (idx) at (3.3,-8.2)
  {\textbf{SQLite index}\\ one row per program; \textsc{vulnerable} / \textsc{safe}; \textsc{unknown} excluded};
\node[box] (samp) at (6.5,-8.2)
  {\textbf{Stratified sample}\\ label $\times$ family, fixed seed, weight $N_h/n_h$};
\node[llm, text width=3.1cm] (pred) at (9.9,-8.2)
  {\textbf{Predictor (source only)}\\ \texttt{tools\_only}: \texttt{cppcheck} + syntax-tree checks;
   \texttt{llm\_only}, \texttt{hybrid}, \texttt{hybrid\_verify}, \texttt{hybrid\_rsi}};
\node[data, text width=3.4cm] (eval) at (13.7,-8.2)
  {\textbf{Evaluation}\\ strict and population-weighted metrics, bootstrap 95\% CI, $\pm$2-line localization};
% ---------------- Arrows ----------------
\draw[flow] (in) -- (rules);
\draw[flow] (rules) -- (sand);
\draw[flow] (sand) -- (suite);
\draw[flow] (suite) -- (an);
\draw[flow] (an) -- node[right, font=\tiny, align=left] {results,\\ raw logs} (led);
\draw[flow] (led) -- (rev);
\draw[flow] (rev) -- (score);
\draw[flow] (score) -- (rep);
\draw[aux] (model.north) -- (rev.south);
\draw[aux] (model.south) -- (pred.north);
\node[font=\tiny, align=left, anchor=west] at (model.east) {also used by\\ test generation};
\draw[flow] (fa) -- (idx);
\draw[flow] (idx) -- (samp);
\draw[flow] (samp) -- node[above, font=\tiny] {source} (pred);
\draw[flow] (pred) -- (eval);
\coordinate (lowY) at (0,-9.55);
\draw[aux] (idx.south) -- (idx.south |- lowY)
  -| node[pos=0.25, below, font=\tiny] {labels, read only after all predictions are written} (eval.south);
% ---------------- Lanes ----------------
\coordinate (laneAtop) at (0,1.6);
\coordinate (laneBtop) at (0,-6.75);
\coordinate (laneBbot) at (0,-9.95);
\begin{scope}[on background layer]
  \node[lane, fit=(in)(an)(rev)(rep)(laneAtop)] (laneA) {};
  \node[lanelabel] at (laneA.north west) {Grading mode};
  \node[lane, fit=(fa)(eval)(pred)(laneBtop)(laneBbot)] (laneB) {};
  \node[lanelabel] at (laneB.north west) {Benchmark mode (FormAI-v2)};
\end{scope}
\end{tikzpicture}%
}
\caption{Workflow of the grading pipeline (top) and the FormAI benchmark (bottom).
Blue boxes are deterministic, orange boxes call the local LLM, and grey boxes are
data. Dashed arrows show the shared model and the labels, which are withheld until
every prediction has been written.}
\label{fig:flow}
\end{figure*}

\subsection{Local Model and Review Strategy}
The LLM is a local Qwen3.5-2B model served by \texttt{llama-server} with an
8192-token context. Every structured stage uses the server's
\texttt{response\_format} JSON-schema mode, so the decoder can only emit
schema-valid JSON. Requests use temperature 0, seed 42, and
\texttt{cache\_prompt=false}, because reusing the KV cache changes batch
boundaries and can change the output even at temperature 0. The application
checks the server build before starting it (and \texttt{make check} runs the same
test), because JSON-schema requests with Qwen
templates need llama.cpp build b12005 or later. Source code and logs are always
marked as data, not instructions.

The model has three roles. It describes what the program does, finds defects
during RSI, and rates metrics that no analyzer measured. It never sets the score
of a metric that an executed test or a completed authoritative analyzer measured,
and every finding must cite evidence IDs that exist in the ledger.

% =========================================================
% III. IMPLEMENTATION WORKFLOW
% =========================================================
\section{Detailed Implementation Workflow}

\subsection{Grading Rules}
A \texttt{rules.json} file contains the assignment statement
(\texttt{problem}), a \texttt{problem\_contract} (input ranges, operand types,
output rule), time and memory limits, the submission path, an optional
\texttt{reference\_solution}, an optional \texttt{test\_generation} block, and a
list of metrics. Each metric names an analyzer, a point cap, and analyzer
parameters. The loader resolves paths relative to the project root, keeps paths
from escaping it, and validates every parameter (timeouts, thresholds, compiler
flags, \texttt{cppcheck} severities, test-file structure and unique test IDs). It
sets each metric's weight to its point cap divided by 100 and rejects a
configuration whose graded, enabled metrics do not total 100 points (weights
summing to 1.0). The
repository includes nine teaching assignments (memory errors, a leap-year logic
error, integer-division averaging, brute-force triplets, a buffer overflow, and
four dynamic-programming tasks) and five LeetCode-style problems with reference
solutions. Each comes with a deliberately flawed submission and the defects the
grader is expected to find.

\subsection{Docker Sandbox}
\texttt{DockerSandbox} starts a uniquely named container from
\texttt{agent-sandbox:latest} as UID/GID~10001, with no network, all Linux
capabilities dropped, \texttt{no-new-privileges}, a read-only root filesystem, a
64-process limit, one CPU, and the memory limit from \texttt{rules.json}. Only the
\path{/workspace} (256\,MB) and \path{/tmp} (64\,MB) tmpfs mounts are writable.
Every command has a timeout (at most 120\,s), output is truncated at 64\,KB, and
uploaded files are limited to 2\,MB. The container is removed after the run. If
Docker is unavailable, the source-only analyzers and the LLM review still run, and
the execution metrics are reported as unavailable.

\subsection{Preflight}
Before the analyzers run, the preflight step compiles the submission once
(\texttt{-std=c11 -Wall -Wextra -O2}) and runs it on the first test input. If it
reports a compile failure, a compile timeout, or a runtime failure, the
execution-dependent analyzers are skipped and marked inconclusive instead of
failed.

\subsection{Test Generation and Execution Oracles}
\label{sec:testgen}
Instructor tests are JSON arrays of \texttt{\{id, input, expected\}}. When
\texttt{test\_generation} is enabled, \path{test_generation.py} builds additional
cases:
\begin{enumerate}
    \item One LLM call analyzes the submission to suggest test dimensions. The \texttt{pycparser} syntax tree adds line-anchored boundary targets, for example \texttt{a < b} at line 31 leads to a test with \texttt{a == b}. Neither source decides what is correct.
    \item The model proposes inputs in small batches (up to 3 cases in the first response and 4 in each later one, at most 500 characters per input), so a response cannot run past the token limit. For assignments that state an exact input shape (an $R \times C$ grid, or $n$ followed by $n$ values), deterministic structural variants are added within the stated bounds.
    \item Expected outputs come from an \emph{execution oracle}, in this order of authority: the instructor's reference solution compiled and run in the sandbox; otherwise a Python reference program that the model writes from the assignment alone (it never sees the student's code), trusted only if it reproduces every instructor case and otherwise required to agree with a second, independent model call; otherwise, without a sandbox, agreement between two model calls, marked as not execution-checked.
    \item Malformed, duplicate, out-of-contract, and unverified candidates are discarded. If no case of the first batch survives, one bounded refill request is made.
\end{enumerate}
The verified suite is written to \path{.generated_tests/} under a SHA-256
fingerprint of the assignment, contract, model, seed, temperature, requested case
count, coverage focus, submission, generator version, instructor cases, and
reference solution. A later run with the same fingerprint reuses the
frozen suite, a conflicting suite is never silently overwritten, and
\texttt{require\_frozen\_suite} makes grading fail rather than regenerate. In the
interface, generation is a separate step so the instructor can review the cases
before grading.

\subsection{Analyzers and Metrics}
Table~\ref{tab:metrics} lists the nine default metrics. \emph{I/O tests} and the
\emph{edge-case runner} compile with \texttt{-std=c11 -Wall -Wextra -O2} and
compare the trimmed actual and expected output per test, recording exit code,
timeout, and truncation. Possibly truncated output never counts as a pass.
\emph{Compiler diagnostics} (\texttt{-std=c11 -Wall -Wextra -Wpedantic}) score 0
if compilation fails and otherwise deduct 5 points per warning. \emph{Static analysis}
runs \texttt{cppcheck} with per-severity penalties (error 25, warning 10,
performance 4, style 3, portability 3, information 0.5). Messages about
\texttt{cppcheck}'s own analysis (such as \texttt{missingIncludeSystem} or
\texttt{checkersReport}) are kept as tool notes and are not counted as defects.
\emph{Modularity}, \emph{complexity}, and \emph{readability} use heuristics:
function count and helper functions, cyclomatic complexity and function length
(via \texttt{lizard}), lines over 100 characters, single-character identifiers,
and the absence of comments. \emph{Similarity} compares 8-token shingles after removing
comments and literals and replacing identifiers. Credit is full at a maximum
similarity of 0.35 or less and falls linearly to zero at 0.85, and the full
pairwise matrix of the cohort is reported. \emph{Runtime profiling} records
repeated wall-clock samples and per-test medians.

\begin{table}[!t]
\centering
\caption{Default Metrics (Total 100 Points) and Score Authority}
\label{tab:metrics}
\footnotesize
\rowcolors{2}{rowgray}{white}
\begin{tabularx}{\columnwidth}{@{}L{1.55cm} L{2.05cm} C{0.6cm} Y@{}}
\toprule
\rowcolor{headgray}
\textbf{Metric} & \textbf{Analyzer} & \textbf{Pts} & \textbf{Score source} \\
\midrule
correctness & \texttt{io\_tests} & 20 & Test pass rate \\
logical & \texttt{edge\_case\_runner} & 10 & Test pass rate \\
syntax & \texttt{compiler\_diagnostics} & 10 & Analyzer \\
memory\_safety & \texttt{memory\_safety} & 10 & Analyzer \\
similarity & \texttt{similarity\_flag} & 10 & Analyzer \\
structure & \texttt{static\_analysis} & 10 & LLM level; analyzer fallback \\
modularity & \texttt{modularity} & 10 & LLM level; analyzer fallback \\
complexity & \texttt{complexity\_estimator} & 10 & LLM level; analyzer fallback \\
readability & \texttt{readability} & 10 & LLM level; analyzer fallback \\
\bottomrule
\end{tabularx}
\end{table}

\subsection{Memory-Safety Analysis}
The memory-safety analyzer first runs GCC's static analyzer
(\texttt{-fanalyzer}) on the source. It then builds a debug binary
(\texttt{-g -O0}) and runs it under Valgrind Memcheck (\texttt{--leak-check=full},
\texttt{--error-exitcode=99}) for each test input, with a separate log per run. It
also builds a second binary with \texttt{-fsanitize=address,undefined} and runs it
on the same inputs. Each tool gets the fraction of runs that were clean. The
metric score is the minimum of the two when both complete, and the Valgrind score
alone when the sanitizer runs are inconclusive. A missing log, a timeout, a truncated log, or a missing
\texttt{ERROR SUMMARY} makes the Valgrind result inconclusive, never clean. These
checks cover only the inputs that were executed.

\subsection{Syntax-Tree Facts and Safety Properties}
\path{c_ast.py} parses the submission with \texttt{pycparser}. It blanks
\texttt{\#include} lines while keeping line numbers, substitutes simple
\texttt{\#define} constants, and declares common library typedefs. It then
extracts functions, recursion, loops and nesting depth, comparisons, I/O calls
(including unchecked \texttt{scanf} results), allocations and frees, reads after
\texttt{free}, divisions, narrow accumulators, and fixed-size arrays.
\path{c_properties.py} uses these facts for line-anchored checks of the property
classes that ESBMC verifies: unbounded \texttt{\%s}/\texttt{gets} input, unchecked
allocation results, loop bounds that reach an array's size, unchecked
input-derived indices, divisors, and arithmetic, leaks, and use after free. These
checks are flow-insensitive approximations and are treated as evidence, not
proofs.

\subsection{Evidence Ledger}
\path{evidence.py} builds the ledger from the rules, the analyzer output, and the
source code only. Each fact gets a stable ID: a failed test with its input,
expected and actual output, a sanitizer or Valgrind error, a compiler warning, a
syntax-tree fact, or a source line range. Observed failures and defects are listed
first. Model text is never added to the ledger as evidence. Cited IDs that are
not in the ledger are dropped, and a finding's line range must lie within the
source lines the model was shown. Long prose sections
can be shortened by the same model when they exceed a size threshold, but source
code, line numbers, finding identities, and measurements are never compressed. If
compression fails validation, a deterministic compaction is used instead.

\subsection{LLM Review: Understanding, RSI, and Adjudication}
\path{reviewer.py} runs the review in stages:
\begin{enumerate}
    \item \textbf{Program understanding:} one request describes what the source actually does, what the assignment expects, how the two align, the algorithm, and the functions. If it fails, a deterministic summary of the source is used instead.
    \item \textbf{RSI passes (at most 3):} the first pass reviews the whole program against the contract and must explain each observed failure in the ledger. Each later pass receives the previous complete JSON response and the defect evidence that no finding cites yet, and must keep, correct, remove, or add findings. The loop stops early when a pass reproduces the previous findings exactly with nothing left to reconcile, a fixed point at temperature~0.
    \item \textbf{Deterministic prevalidation:} after each pass, Python removes configuration complaints, self-contradictions, claims disproved by the source (for example a ``missing input check'' where \texttt{scanf}'s result is checked), memory-safety claims with no pointer or indexed access at the cited line, optional style preferences, generic non-falsifiable claims, and duplicates before they reach the next pass.
    \item \textbf{Adjudication:} one separate request checks each remaining candidate for source location, relevance to the contract, logical consistency, and evidence. Rejected candidates are recorded in the report and do not affect any score.
    \item \textbf{Metric rating:} one request rates only the metrics without a deterministic score, as an integer level from 0 to 10 using rubric anchors (10 = fully meets; 0 = not met, with a specific failure). Python converts each level to points.
\end{enumerate}
The adjudication and rating stages allow up to three bounded repair attempts for
malformed JSON.
If the server cannot start, a fallback evaluation that is clearly marked as such
is built from the evidence.

\subsection{Deterministic Scoring}
\path{scoring/overall.py} selects one score per metric. For \texttt{correctness}
and \texttt{logical}, the observed pass rate of the executed tests is the score,
whatever the model says. For the other analyzers in the first group of
Table~\ref{tab:metrics}, the completed analyzer's score is used. For the remaining
metrics, a valid model level is used, and the analyzer's score is the fallback.
A model rating whose rationale says the metric could not be measured is
discarded, and a model rating can never replace missing execution evidence. Each
metric records its origin (\texttt{deterministic\_test\_results},
\texttt{deterministic\_analyzer}, \texttt{llm}, or \texttt{unmeasured}), its
available and missing evidence, and a 5-point quality band.

The grade is the sum of the awarded points. If any graded metric is unmeasured,
no grade is calculated: the status is \texttt{failed\_incomplete} and the missing
metrics are listed with their causes. If a grade is calculated but an LLM stage
used a fallback or rejected candidates, the status is
\texttt{completed\_with\_warnings}.

\subsection{Reporting and Reproducibility}
Each run writes a JSON report and a text report. The JSON report contains the
source SHA-256, the rules SHA-256, tool versions (GCC, \texttt{cppcheck},
Valgrind, \texttt{lizard}), the runtime configuration (model, temperature, RSI
iterations, token budgets), every analyzer result with raw logs, the evidence
ledger, the RSI history, accepted and rejected findings, the score provenance, and
timing for each stage. The text report gives the final grade, each metric's
points and source with a short rationale, the validated findings with line
references, the test-suite origin and oracle, the program understanding, and the
stage timings. Values that are not
valid JSON (such as non-finite floats) are converted before writing, so a
completed run always produces a report.

\subsection{FormAI Benchmark Pipeline}
\path{formai_benchmark.py} runs one benchmark at a time. It draws a stratified
random sample (strata are label $\times$ vulnerability family, or label $\times$
generating model) with balanced or proportional allocation and a fixed seed.
Every sampled program carries the design weight $N_h/n_h$. The predictor receives
the source code only, and the labels are read after all predictions are written.
Five prediction modes are implemented:
\begin{itemize}
    \item \texttt{tools\_only}: \texttt{cppcheck} plus the syntax-tree property checks, merged to one finding per (line, family);
    \item \texttt{llm\_only}: an LLM source audit with exact source excerpts;
    \item \texttt{hybrid}: tools plus the LLM audit;
    \item \texttt{hybrid\_verify}: tools plus LLM claims that pass a separate verification;
    \item \texttt{hybrid\_rsi}: tools plus the LLM audit, with RSI adjudicating every LLM-originated candidate against source and analyzer evidence. Confirmed tool positives are kept.
\end{itemize}
A stage that fails or an unresolved positive candidate makes the prediction
inconclusive, not safe. Each run folder under \path{reports/formai/} holds
\path{config.json}, \path{sample.json}, \path{inputs.jsonl} (source only),
\path{predictions.jsonl}, \path{labels.json}, \path{evaluation.json}, and
\path{report.md}.

% =========================================================
% IV. EVALUATION PROTOCOL AND METRICS
% =========================================================
\section{Evaluation Protocol and Metrics}

\subsection{Evaluation Protocol}
Ground truth is defined per program: \textsc{vulnerable} or \textsc{safe}. Headline
metrics are \emph{strict}: an inconclusive prediction counts as an error, that is,
a missed vulnerability on a vulnerable program and a false alarm on a safe one.
Metrics over decided programs only are reported next to them. Because the sample
is stratified, we report two sets of figures. The \emph{sample} figures describe
the drawn programs. The \emph{dataset estimates} weight each program by $N_h/n_h$
so that they describe the whole eligible population, whose prevalence (88\%
vulnerable) differs from a balanced sample's. Precision in particular depends on
prevalence. 95\% confidence intervals come from a stratified bootstrap (resampling
within each stratum, fixed seed, 1000 resamples for the larger runs). A ratio with
a zero denominator is reported as undefined, never as zero.

\subsection{Metrics for Comparison}
With TP, FP, FN, and TN denoting true positives, false positives, false
negatives, and true negatives:
\begin{align}
\text{Precision} &= \frac{TP}{TP+FP}, \qquad
\text{Recall} = \frac{TP}{TP+FN}, \\
F_1 &= 2\,\frac{\text{Precision}\cdot\text{Recall}}{\text{Precision}+\text{Recall}}.
\end{align}
We also report specificity, balanced accuracy, accuracy, and the Matthews
correlation coefficient (MCC). For localization, a finding \emph{hits} a
vulnerability when its line is within $\pm 2$ lines of an ESBMC vulnerable line;
\emph{same family} also requires the vulnerability family to match. ESBMC checks
a bounded program, so a finding away from its lines is not necessarily wrong, and
location precision is a lower bound. Detection and location are also reported per
family and per stratum, together with the runtime per program.

% =========================================================
% V. RESULTS
% =========================================================
\section{Results}

\subsection{Static-Tool Detection on FormAI-v2}
Table~\ref{tab:formai_tools} shows the \texttt{tools\_only} mode on a balanced
sample of 600 programs (300 vulnerable, 300 safe, seed 7). Every program was
decided. The detector is conservative: it raises few false alarms (13 of 300 safe
programs) but finds about two thirds of the vulnerable programs. Among the 190
detected vulnerable programs, the primary finding was on an ESBMC line in 80.5\%
of cases, and the line and family both matched in 86.3\% of cases for at least one
finding. A 40-program run with a different seed gave lower recall (50.0\%)
within wide confidence intervals. The static stage takes about 0.14\,s per
program.

\begin{table}[!t]
\centering
\caption{FormAI-v2, \texttt{tools\_only}, 600 Programs (Strict)}
\label{tab:formai_tools}
\footnotesize
\rowcolors{2}{rowgray}{white}
\begin{tabularx}{\columnwidth}{@{}L{2.4cm} C{1.2cm} Y@{}}
\toprule
\rowcolor{headgray}
\textbf{Metric} & \textbf{Sample} & \textbf{Dataset estimate (95\% CI)} \\
\midrule
Recall & 63.3\% & 63.5\% (58.5--68.5) \\
Specificity & 95.7\% & 95.7\% (93.3--98.0) \\
Precision & 93.6\% & 99.1\% (98.6--99.6) \\
$F_1$ & 75.5\% & 77.4\% (73.7--81.1) \\
Balanced accuracy & 79.5\% & 79.6\% (76.9--82.4) \\
MCC & 0.624 & 0.384 (0.347--0.425) \\
\midrule
\multicolumn{3}{@{}l}{Confusion: TP 190, FN 110, FP 13, TN 287} \\
\bottomrule
\end{tabularx}
\end{table}

Table~\ref{tab:family} shows detection by family. NULL dereferences (mostly
unchecked \texttt{malloc} results) are found and located reliably. Out-of-bounds
and invalid-pointer vulnerabilities are often detected in the right program but
rarely at ESBMC's line, because the syntax-tree checks do not track pointer
aliasing or indices computed at run time.

\begin{table}[!t]
\centering
\caption{\texttt{tools\_only} Detection by ESBMC Family (600-Program Run)}
\label{tab:family}
\footnotesize
\rowcolors{2}{rowgray}{white}
\begin{tabularx}{\columnwidth}{@{}Y C{0.9cm} C{1.2cm} C{1.4cm}@{}}
\toprule
\rowcolor{headgray}
\textbf{Family} & \textbf{Progs} & \textbf{Detected} & \textbf{Located (family)} \\
\midrule
Input buffer overflow & 140 & 61.4\% & 49.3\% \\
NULL dereference & 119 & 87.4\% & 83.2\% \\
Out of bounds & 39 & 46.2\% & 2.6\% \\
Invalid pointer & 42 & 50.0\% & 0.0\% \\
Memory leak & 25 & 76.0\% & 20.0\% \\
Integer overflow & 28 & 67.9\% & 17.9\% \\
Float overflow & 10 & 20.0\% & 0.0\% \\
Division by zero & 3 & 33.3\% & 33.3\% \\
\bottomrule
\end{tabularx}
\end{table}

\subsection{LLM-Assisted Modes}
Table~\ref{tab:formai_modes} compares the LLM-assisted modes. They are much slower
(12--94\,s per program), so they were run on small samples, and their confidence
intervals are wide. On the same 12 balanced programs, \texttt{llm\_only} and
\texttt{hybrid} found every vulnerable program, but \texttt{llm\_only} flagged 4
of the 6 safe programs and \texttt{hybrid} flagged 4 and left 2 undecided: the 2B
model rarely concludes that a program is safe. \texttt{hybrid\_rsi} removed
unsupported positives but left 8 of 12 programs inconclusive, which the strict
metrics count as errors. On a 20-program proportional sample (18 vulnerable),
\texttt{hybrid\_rsi} decided 11 programs, all correctly, and 90.5\% of its findings
fell on ESBMC lines. Its strict recall (55.6\%) is limited by the 9 inconclusive
programs. Specificity estimates from samples with 2--6 safe programs are not
reliable. These results suggest that RSI adjudication improves precision and
localization, but that its inconclusive rate and cost must come down before it can
replace the static detector.

\begin{table}[!t]
\centering
\caption{LLM-Assisted Modes on Small FormAI-v2 Samples (Sample Figures, Strict)}
\label{tab:formai_modes}
\footnotesize
\rowcolors{2}{rowgray}{white}
\begin{tabularx}{\columnwidth}{@{}L{1.55cm} C{0.55cm} C{0.95cm} C{0.65cm} C{0.65cm} C{0.65cm} Y@{}}
\toprule
\rowcolor{headgray}
\textbf{Mode} & \textbf{$n$} & \textbf{Decided} & \textbf{Rec.} & \textbf{Spec.} & \textbf{Prec.} & \textbf{s/prog} \\
\midrule
\texttt{llm\_only} & 12 & 12 & 100\% & 33.3\% & 60.0\% & 12.6 \\
\texttt{hybrid} & 12 & 10 & 100\% & 0.0\% & 50.0\% & 12.1 \\
\texttt{hybrid\_rsi} & 12 & 4 & 50.0\% & 0.0\% & 33.3\% & 83.7 \\
\texttt{hybrid\_rsi}$^{\dagger}$ & 20 & 11 & 55.6\% & 50.0\% & 90.9\% & 94.4 \\
\bottomrule
\multicolumn{7}{@{}l}{\scriptsize $^{\dagger}$Proportional allocation, seed 42 (18 vulnerable, 2 safe); other rows balanced, seed 7.}
\end{tabularx}
\end{table}

\subsection{Grading Runs}
The grading pipeline was run on the example assignments through the Gradio
interface; the reports are kept in \path{reports/gradio/}. A review makes 6--8 LLM
requests (understanding, RSI passes, adjudication, and rating, including any
repair attempts), and grading takes 84--150\,s. Test generation is a separate step
in the interface and is not included in these times. RSI stopped after two of the
three allowed passes in most runs.

One array-processing submission was graded repeatedly during development, so the
code changed between some runs. With the instructor's five correctness tests, four
consecutive runs gave the same total (75.00/100), and a later run, after changes
to the reviewer, gave 86.00. In all five runs the measured metrics (correctness 4/5
tests, logical, syntax, similarity, and memory safety) gave the same 51 points. The
whole difference came from the model's ratings of complexity, structure, and
readability. With generated suites of 39 and 49 cases, the totals were 76.52,
76.52, 78.14, and 85.14. In every run, 60 of the 100 points came from executed
tests or completed analyzers, and the remaining 40 from model ratings of
structure, modularity, complexity, and readability. These model-rated metrics are
the main source of variation between runs and between code versions. The repository contains 371 unit and regression test functions
covering the analyzers, scoring precedence, the evidence ledger, RSI, test
generation, oracles, fallbacks, and the benchmark.

% =========================================================
% VI. IMPLEMENTATION DETAILS
% =========================================================
\section{Implementation Details}

\subsection{Implementation Status}
Table~\ref{tab:plan} summarizes the implemented milestones and where each one is
in the repository.

\begin{table*}[!t]
\centering
\caption{Implemented Milestones}
\label{tab:plan}
\footnotesize
\rowcolors{2}{rowgray}{white}
\begin{tabularx}{\textwidth}{@{}C{0.9cm} L{3.2cm} Y L{3.6cm}@{}}
\toprule
\rowcolor{headgray}
\textbf{Stage} & \textbf{Milestone} & \textbf{Implemented activities} & \textbf{Module} \\
\midrule
1 & Sandbox &
Unprivileged, network-less container with read-only root, dropped capabilities, CPU/memory/PID limits, tmpfs workspace, timeouts, and output caps; graceful degradation without Docker. &
\path{sandbox/docker_sandbox.py} \\
2 & Rule loader &
Validation of the assignment, contract, limits, metrics, weights, point caps, and test paths. &
\path{config/rule_loader.py} \\
3 & Preflight and analyzers &
Compile and smoke check; ten registered analyzers covering tests, compiler, \texttt{cppcheck}, Valgrind, ASan/UBSan, runtime, modularity, complexity, readability, and similarity. &
\path{preflight.py}, \path{analyzers/} \\
4 & Syntax-tree analysis &
\texttt{pycparser} facts, edge-case targets, and ESBMC-class safety-property checks. &
\path{c_ast.py}, \path{c_properties.py} \\
5 & Test generation &
Batched input generation, structural variants, execution oracles (instructor or model-written reference), independent verification, and frozen fingerprinted suites. &
\path{testing/} \\
6 & LLM review &
Evidence ledger, program understanding, bounded RSI with convergence stop, prevalidation, adjudication, rubric levels, and context compression. &
\path{llm/} \\
7 & Deterministic scoring &
Evidence-priority metric selection, point sum, unmeasured-metric policy, and score provenance. &
\path{scoring/} \\
8 & Reporting and UI &
JSON and text reports with hashes, tool versions, and timings; Gradio grading and benchmark tabs. &
\path{reporting/}, \path{ui/} \\
9 & FormAI benchmark &
Dataset download and indexing, stratified sampling with design weights, five prediction modes, strict and weighted metrics, bootstrap CIs, and localization. &
\path{benchmark/} \\
\bottomrule
\end{tabularx}
\end{table*}

\subsection{Development Commands}
On Linux, \texttt{make} installs Python dependencies with \texttt{uv}, Docker,
\texttt{llama-server}, the sandbox image, and the model, then opens the interface
at \url{http://127.0.0.1:7860}:
\begin{lstlisting}
make            # setup + launch UI
make sandbox    # rebuild Dockerfile.sandbox
make check      # verify llama-server build
make test       # run the test suite
\end{lstlisting}
Grading and benchmarking from the command line:
\begin{lstlisting}
python grade.py --rules path/to/rules.json \
  [--source file.c] [--refresh-suite]
PYTHONPATH=src python -m \
  grading_agent.benchmark.formai_benchmark \
  --dataset v2 -n 50 --mode hybrid_rsi
\end{lstlisting}
Settings (model, ports, context size, RSI iterations, sandbox limits, analyzer
thresholds) are read from environment variables or \path{.env} in
\path{config/settings.py}. The interface can override a whitelisted subset for
each request without affecting concurrent jobs.

% =========================================================
% VII. WORK DISTRIBUTION
% =========================================================
\section{Distribution of Work Among Group Members}

The group has four members; the allocation below is provisional and should be
updated to reflect actual contributions.

\begin{table}[!t]
\centering
\caption{Work Allocation}
\label{tab:work}
\footnotesize
\rowcolors{2}{rowgray}{white}
\begin{tabularx}{\columnwidth}{@{}L{2.5cm} L{1.7cm} Y@{}}
\toprule
\rowcolor{headgray}
\textbf{Member} & \textbf{Roll No.} & \textbf{Responsibility} \\
\midrule
Gaurav Chandel & 2611MC17 & AI Integration \\
Prajodeep Patil & 2611MC11 & AI Integration \\
Aditya Sharma & 2611CS11 & Docker Implementation \\
Anshul Kaushal & 2611AI10 & Docker Implementation \\
\bottomrule
\end{tabularx}
\end{table}

The functional areas are: (i) input handling and rule validation (grading schema,
JSON loading, configuration validation); (ii) Docker sandbox and execution
(compiler, controlled execution, sanitizers, failure recording); (iii) analysis
and scoring (analyzers, evidence ledger, deterministic grading); and (iv) LLM
integration, test generation, benchmarking, and reporting.

% =========================================================
% VIII. USEFULNESS AND LIMITATIONS
% =========================================================
\section{Usefulness and Limitations}

\subsection{Potential Real-World Usefulness}
\begin{enumerate}
    \item \textbf{Coursework at scale.} An instructor can run the same frozen test suite and analyzers against every submission and get consistent evidence for each test.
    \item \textbf{Actionable feedback.} A failing test reports the input, expected and actual output, exit code, and timeout status. Valgrind and sanitizer logs and line-anchored findings point to the responsible statements.
    \item \textbf{Auditable grades.} Each metric's points carry their origin and evidence, and the source hash, rules hash, and tool versions are recorded with the grade.
    \item \textbf{Fewer hand-written tests.} Generated tests with execution-checked expected outputs extend a few instructor examples to boundary cases.
    \item \textbf{Triage for human review.} Similarity flags and inconclusive metrics identify submissions that need a person to look at them. A similarity score alone is not proof of copying.
\end{enumerate}

\subsection{Concrete Limitations and Failure Cases}
\begin{enumerate}
    \item \textbf{Model-rated metrics.} With the default rules, 40 of 100 points come from the 2B model's 0--10 ratings, which remain the main source of variation between runs. \emph{Mitigation:} rubric anchors, ledger citations, and analyzer fallbacks; instructors can raise the weight of measured metrics.
    \item \textbf{Test coverage.} A program can pass every configured case but fail an untested input, and dynamic tools report only executed paths. \emph{Mitigation:} syntax-tree edge-case targets, structural variants, and frozen suites the instructor can review.
    \item \textbf{Oracle trust.} Without an instructor reference, expected outputs come from a model-written program that is checked against instructor cases or a second model. A wrong reference produces wrong expected outputs. \emph{Mitigation:} provide \texttt{reference\_solution}; the report states which oracle was used.
    \item \textbf{LLM detection quality.} The LLM-assisted benchmark modes have low specificity or high inconclusive rates and cost 12--94\,s per program. \emph{Mitigation:} keep tool positives authoritative and count inconclusive results as errors.
    \item \textbf{Static-check approximations.} The syntax-tree checks are flow-insensitive and miss aliasing, which explains the low location rates for out-of-bounds and invalid-pointer vulnerabilities.
    \item \textbf{Sandbox residual risk.} Isolation reduces but does not remove exposure to kernel or Docker vulnerabilities. \emph{Mitigation:} no secrets or host mounts, an unprivileged user, and an updated runtime.
\end{enumerate}

The system is a grading assistant, not a replacement for instructor judgment. Its
reports separate measured results from model estimates so that a person can check
the estimates.

% =========================================================
% IX. CONCLUSION
% =========================================================
\section{Conclusion}
The Code Checking Assistant grades C submissions end to end. A sandboxed analyzer
pipeline measures correctness, compilation, static diagnostics, and memory safety
with both Valgrind and sanitizers. Test generation produces execution-checked,
frozen suites. A local 2B model, constrained to cite a deterministic evidence
ledger, reviews the program in bounded RSI passes. Python computes the grade with
measured results taking precedence, and refuses to produce a grade when a metric
has no evidence. On FormAI-v2, the static detector reaches 63.3\% recall at
95.7\% specificity on 600 programs, and RSI adjudication gives precise,
well-located LLM findings at a high cost and inconclusive rate.

Future work is to reduce the share and variance of model-rated points, lower the
inconclusive rate and latency of \texttt{hybrid\_rsi}, evaluate it on larger
samples, add flow-sensitive checks for pointers and indices, and measure agreement
with human graders on real course submissions.

% =========================================================
% REFERENCES
% =========================================================
\balance
\begin{thebibliography}{1}
\bibitem{formai}
N. Tihanyi \emph{et al.}, ``The FormAI dataset: Generative AI in software
security through the lens of formal verification,'' in \emph{Proc. PROMISE},
2023. DOI: \url{https://doi.org/10.1145/3617555.3617874}.

\bibitem{formaiv2}
N. Tihanyi \emph{et al.}, ``Do neutral prompts produce insecure code? FormAI-v2
dataset: Labelling vulnerabilities in code generated by large language models,''
arXiv:2404.18353, 2024.

\bibitem{formairepo}
FormAI Dataset repository. [Online]. Available:
\url{https://github.com/FormAI-Dataset/FormAI-dataset}

\bibitem{esbmc}
M. R. Gadelha, F. R. Monteiro, J. Morse, L. C. Cordeiro, B. Fischer, and
D. A. Nicole, ``ESBMC 5.0: An industrial-strength C model checker,'' in
\emph{Proc. ASE}, 2018, pp.~888--891.
\end{thebibliography}

\end{document}
